\usepackage{fontspec}
\setsansfont{cmunss.otf}[BoldFont=cmunsx.otf,ItalicFont=cmunsi.otf]
\usepackage{polyglossia}
\setdefaultlanguage{russian}
\setotherlanguage{english}
\newfontfamily\cyrillicfontsf{cmunss.otf}[BoldFont=cmunsx.otf,ItalicFont=cmunsi.otf]
\usepackage{booktabs,tabularx,amsmath}
\usepackage{tikz}
\usetikzlibrary{arrows.meta,positioning,calc}
\usepackage{array}
\definecolor{Ink}{HTML}{182636}
\definecolor{Muted}{HTML}{495D72}
\definecolor{Route}{HTML}{126A63}
\definecolor{Trace}{HTML}{2D5C86}
\definecolor{Rule}{HTML}{CCD8DF}
\definecolor{Paper}{HTML}{FFFFFF}
\setbeamercolor{normal text}{fg=Ink,bg=Paper}
\setbeamercolor{structure}{fg=Route}
\setbeamercolor{frametitle}{fg=Ink}
\setbeamercolor{footline}{fg=Muted}
\setbeamerfont{frametitle}{size=\fontsize{17}{20}\selectfont,series=\bfseries}
\setbeamersize{text margin left=0.75cm,text margin right=0.75cm}
\setbeamertemplate{navigation symbols}{}
\setbeamertemplate{frametitle}{\vspace{0.28cm}\insertframetitle\par\vspace{0.18cm}}
\setbeamertemplate{footline}{\hfill\usebeamercolor[fg]{footline}\fontsize{7}{8}\selectfont\insertframenumber/\inserttotalframenumber\hspace*{0.75cm}\vspace{0.28cm}}
\hypersetup{colorlinks=true,urlcolor=Trace}
\newcommand{\highlight}[1]{{\color{Route}\bfseries #1}}
\newcommand{\muted}[1]{{\color{Muted}#1}}
\newcommand{\slidebody}{\fontsize{10.8}{13.8}\selectfont}
\newcommand{\slidenote}[1]{\par\vspace{0.24cm}{\color{Muted}\fontsize{8.2}{10.5}\selectfont #1\par}}
\newcommand{\legendmark}[1]{\tikz[baseline=-0.6ex]\fill[#1] (0,0) rectangle (0.18,0.12);}
\newcommand{\ce}{\operatorname{CE}}
\newcommand{\relu}{\operatorname{BiReLU}}
\newcommand{\projecttitle}{Компиляция статической маршрутизации\par с нелинейностью до агрегации}
\newcommand{\PresenterName}{}
\newcommand{\PresenterAffiliation}{}
\newcommand{\presenter}{\ifx\PresenterName\empty\else\vspace{0.45cm}\muted{\PresenterName\ifx\PresenterAffiliation\empty\else\par\PresenterAffiliation\fi}\fi}

\newcommand{\maindiagram}{%
\begin{tikzpicture}[x=1cm,y=1cm,
  edge/.style={draw=Muted,-{Latex[length=1.8mm]},line width=0.52pt},
  token/.style={font=\fontsize{9.5}{12}\selectfont,align=center,text=Ink},
  routepart/.style={token,text=Route},
  exactpart/.style={token,text=Trace}]
  \node[token] (x) at (0,0) {Входы\\$x_1,\ldots,x_N$};
  \node[routepart] (stats) at (3.15,0.74) {Общие статистики\\по интервалам};
  \node[exactpart] (select) at (3.15,-0.74) {Входы явных\\маршрутов $S_m$};
  \node[routepart] (tail) at (6.85,0.74) {Сумма BiReLU\\по порогу $-b_m/w_m$};
  \node[exactpart] (correction) at (6.85,-0.74) {Поправка для\\явных маршрутов};
  \node[token] (out) at (10.2,0) {Выход\\$h_m^Q(x)$};
  \draw[edge] (x.east) -- ++(0.33,0) |- (stats.west);
  \draw[edge] (x.east) -- ++(0.33,0) |- (select.west);
  \draw[edge,Route] (stats.east) -- (tail.west);
  \draw[edge,Trace] (select.east) -- (correction.west);
  \draw[edge] (tail.east) -- ++(0.3,0) |- (out.west);
  \draw[edge] (correction.east) -- ++(0.3,0) |- (out.west);
\end{tikzpicture}}

\newcommand{\oneslidecontent}{%
\slidebody
Слой усредняет отклики BiReLU по обученным вероятностям маршрутов.
При выводе компилятор хранит явные маршруты и остаточную массу групп.
\par\vspace{0.4cm}
\centering\maindiagram\par
\vspace{0.46cm}
\begin{columns}[T,onlytextwidth]
\begin{column}{0.31\textwidth}
{\color{Route}\fontsize{25}{28}\selectfont\bfseries $2{,}81\times$}\par
быстрее точного кеша
\end{column}
\begin{column}{0.31\textwidth}
{\color{Route}\fontsize{25}{28}\selectfont\bfseries $10{,}3\times$}\par
меньше явной памяти
\end{column}
\begin{column}{0.31\textwidth}
{\color{Route}\fontsize{25}{28}\selectfont\bfseries $-0{,}35$}\par
п.п. точности на тесте
\end{column}
\end{columns}
\raggedright
\slidenote{База сравнения — исходная маршрутизирующая сеть с точным кешем.
CIFAR-10, точность по трём запускам. Замеры Apple M3 Pro CPU, FP32, batch 1, seed 0.}
}
